% **************************************************************************************************************
%
%  Copyright 2020-2026 Robert Bosch GmbH
%
%  Licensed under the Apache License, Version 2.0 (the "License");
%  you may not use this file except in compliance with the License.
%  You may obtain a copy of the License at
%
%      http://www.apache.org/licenses/LICENSE-2.0
%
%  Unless required by applicable law or agreed to in writing, software
%  distributed under the License is distributed on an "AS IS" BASIS,
%  WITHOUT WARRANTIES OR CONDITIONS OF ANY KIND, either express or implied.
%  See the License for the specific language governing permissions and
%  limitations under the License.
%
% **************************************************************************************************************

\section{Introduction}

\textbf{ProcessHub} is a Python package for process lifecycle management with ZMQ-based coordination.

It provides:

\begin{itemize}
   \item \textbf{Process Lifecycle Management}: Start, stop, and restart processes across distributed systems
   \item \textbf{Health Monitoring}: Automatic detection and recovery from crashed processes
   \item \textbf{Restart Coordination}: Multi-client restart orchestration with acknowledgment protocol
   \item \textbf{Pluggable Transport}: ZMQ for production, in-memory for testing
   \item \textbf{Decoupled Architecture}: Core logic separated from I/O for testability
\end{itemize}

\subsection{Target Audience}

This package is designed for developers who need to:

\begin{itemize}
   \item Manage multiple processes in a distributed system
   \item Coordinate process restarts across multiple clients
   \item Monitor process health and handle failures automatically
   \item Build robust process management into their applications
\end{itemize}

\subsection{Prerequisites}

\begin{itemize}
   \item Python 3.10 or higher
   \item psutil package (required)
   \item pyzmq package (optional, for ZMQ transport)
\end{itemize}
